In this proposal, I set out the framing for a paper which estimates the causal effect of new rail station access on nearby residents' labor-market outcomes, focusing on worker flows and income. New stations may improve outcomes through three mechanisms: more reliable commuting to existing jobs, expanded job-search radius across the metro area, and broader referral networks. Identifying these effects is complicated by non-random station placement, residential sorting, and anticipatory adjustment. I address these challenges by exploiting quasi-random variation in station opening timing generated by construction delays, embedding the design in a stacked difference-in-differences framework that yields robust event-study estimates under staggered adoption. A novel dataset of U.S. rail projects since 2000 links planned opening dates — drawn from Draft Environmental Impact Statements — to realized opening dates and station geographies. Labor-market outcomes come from Census LODES, and transit exposure is measured via walking-time isochrones computed using the \texttt{r5r} routing package.